\begin{table}[tbp]
\caption{Questions, evidence, and products in an open source think tank.}
\label{tab:ch10-audiences}
\centering\small\setlength{\tabcolsep}{4pt}\renewcommand{\arraystretch}{1.19}
\begin{tabularx}{\linewidth}{@{}>{\raggedright\arraybackslash}p{.15\linewidth}>{\raggedright\arraybackslash}p{.25\linewidth}>{\raggedright\arraybackslash}X>{\raggedright\arraybackslash}p{.18\linewidth}@{}}
\toprule
Audience & Decision question & Relevant evidence & Product \\
\midrule
National institutions & Where are capabilities and collaborative relationships concentrated? & Participation, contribution, influence, and their attribution rules & Development report \\
\addlinespace
Regional institutions & Which local capabilities merit further investigation or support? & Consistent city series, project composition, roles, and local conditions & Regional profile \\
\addlinespace
Enterprises & How is the project portfolio sustained? & Project scope, contribution structure, review and release continuity & Portfolio dashboard \\
\addlinespace
Foundations and communities & Where do stewardship needs exceed available capacity? & Maintainer roles, continuity, external participation, and governance arrangements & Community assessment \\
\addlinespace
Universities and developers & Where can participation lead to sustained collaborative work? & Project context, contribution pathways, and evidence of learning & Project dossier or educational brief \\
\bottomrule
\end{tabularx}
\tablenote{Adapted from the OpenThink product--audience matrix \citep{openthink}. Rows describe analytical service designs; they do not assert completed engagements with every audience.}
\end{table}
